\section*{अध्याय \ref{ch:g6:water-states} --- पानी अपनी सभी अवस्थाओं में}

\begin{solution}{exo:g6:water-states:1}
\omterm{def:g6:water-states:changes}{गलन} (सुबह घास से पिघलता पाला); \omterm{def:g6:water-states:changes}{ठोसीकरण} (फ़्रीज़र की बर्फ़); \omterm{def:g6:water-states:changes}{वाष्पन}
(सूखते कपड़े, उबलता पतीला); \omterm{def:g6:water-states:changes}{संघनन} (ठंडे \omterm{def:g5:mirrors-reflection:reflection}{दर्पण} पर धुंध); \omterm{def:g6:water-states:changes}{ऊर्ध्वपातन}
(सूखे पाले में सिकुड़ता हिम का ढेर); और सीधे \omterm{def:g3:solids-liquids-gases:gas}{गैस} से \omterm{def:g3:solids-liquids-gases:solid}{ठोस} (शीशे पर पाले
के पंख)।
\end{solution}

\begin{solution}{exo:g6:water-states:2}
ओस: \omterm{def:g6:water-states:changes}{संघनन}। बिना गड्ढा बनाए सिकुड़ता पुतला: \omterm{def:g6:water-states:changes}{ऊर्ध्वपातन} --- ठंडी सूखी
\omterm{def:g2:air-around-us:air}{हवा} में \omterm{def:g3:solids-liquids-gases:solid}{ठोस} सीधे वाष्प बन जाता है। फ़्रीज़र का ``धुआँ'': \omterm{def:g6:water-states:changes}{संघनन} ---
कमरे की वाष्प ठंडी \omterm{def:g2:air-around-us:air}{हवा} में \omterm{def:g2:ice-water-steam:evapcond}{संघनित} होकर नन्हा बादल बना देती है (यानी
बूँदें, \omterm{def:g3:solids-liquids-gases:gas}{गैस} नहीं)।
\end{solution}

\begin{solution}{exo:g6:water-states:3}
तब भी \qty{412}{g}: \omterm{def:g3:measuring-mass:mass}{द्रव्यमान} अवस्था के बदलावों को झेल जाता है ---
बंद बोतल ने न कुछ खोया, न कुछ पाया।
\end{solution}

\begin{solution}{exo:g6:water-states:4}
ढक्कन बंद, ताकि कोई वाष्प अपने ग्रामों समेत उड़ न जाए (हिसाब बंद रहना
चाहिए); और पोंछना इसलिए कि बाहर चिपका पानी या पाला ऐसे बेगाने \omterm{def:g3:measuring-mass:units}{ग्राम}
जोड़ देता जो कमरे के हैं, हमारे हिसाब के नहीं।
\end{solution}

\begin{solution}{exo:g6:water-states:5}
नष्ट नहीं हुआ: \qty{200}{g} पानी वाष्प बनकर रसोई की \omterm{def:g2:air-around-us:air}{हवा} में निकल गया
--- \omterm{def:g6:water-states:changes}{वाष्पन} ने ग्रामों को पतीले से बाहर किया, दुनिया से बाहर नहीं।
\end{solution}

\begin{solution}{exo:g6:water-states:6}
लगभग \qty{1.1}{L} बर्फ़ --- यानी दसवाँ हिस्सा ज़्यादा जगह। \omterm{def:g3:measuring-mass:mass}{द्रव्यमान}
वही रहता है: \qty{1}{kg} भीतर गया, \qty{1}{kg} बाहर आया।
\end{solution}

\begin{solution}{exo:g6:water-states:7}
पाइप में फँसा पानी जमते समय फैलता है, और पाइप काँच की बोतल की तरह ही
हार मान जाता है। माली शरद में नल ख़ाली करा देते हैं ताकि भीतर फैलने को
कुछ बचे ही नहीं।
\end{solution}

\begin{solution}{exo:g6:water-states:8}
बर्फ़ एक \omterm{def:g6:mass-and-volume:volume}{लीटर} का पदार्थ $1.1$ \omterm{def:g6:mass-and-volume:volume}{लीटर} जगह में ठूँस देती है: बराबर जगह के
लिए वह पानी से हल्की है --- और बराबर-जगह-में-हल्का होना ही तैरने की
शर्त है।
\end{solution}

\begin{solution}{exo:g6:water-states:9}
समुद्र (\qty{1}{kg}, \omterm{def:g3:solids-liquids-gases:liquid}{द्रव}) --- \omterm{def:g6:water-states:changes}{वाष्पन} --- वाष्प (\qty{1}{kg})
--- \omterm{def:g6:water-states:changes}{संघनन} --- बादल की बूँदें (\qty{1}{kg}) --- \omterm{def:g6:water-states:changes}{ठोसीकरण} --- हिम
(\qty{1}{kg}) --- \omterm{def:g6:water-states:changes}{गलन} --- पिघला पानी (\qty{1}{kg}) --- नदी से
समुद्र (\qty{1}{kg})। छह पड़ाव, और \omterm{def:g3:measuring-mass:mass}{द्रव्यमान} एक ही, अटल।
\end{solution}

\begin{solution}{exo:g6:water-states:10}
तसल्ली दो। डोलती बर्फ़ पहले ही अपने \emph{\omterm{def:g4:weight-and-mass:weight}{भार}} जितनी शिकंजी हटा चुकी
है; पिघलने पर वह ठीक उतना ही पानी लौटाती है, और साथ ही उस दसवें
हिस्से जितना सिकुड़ भी जाती है जितना वह फूली थी --- यानी वह उसी जगह
में समा जाती है जो टुकड़ा सतह के नीचे पहले से घेरे था। गिलास किनारे तक
भरा रहता है और एक बूँद भी नहीं छलकती।
\end{solution}

\begin{solution}{exo:g6:water-states:11}
\omterm{def:g3:solids-liquids-gases:gas}{गैस} से सीधे \omterm{def:g3:solids-liquids-gases:solid}{ठोस} वाला दरवाज़ा: कमरे की \omterm{def:g2:air-around-us:air}{हवा} की अनदेखी वाष्प जमते शीशे को
छूकर सीधे बर्फ़ बना गई, और \omterm{def:g3:solids-liquids-gases:liquid}{द्रव} अवस्था को छोड़ गई। इसे ``\omterm{def:g2:ice-water-steam:meltfreeze}{जमना}'' इसलिए
नहीं कहते कि शीशे पर कोई \omterm{def:g3:solids-liquids-gases:liquid}{द्रव} कभी था ही नहीं --- पानी सीधे वाष्प
अवस्था से आया।
\end{solution}

\begin{solution}{exo:g6:water-states:12}
तालाब अपनी सतह से जाड़े की \omterm{def:g2:air-around-us:air}{हवा} को \omterm{def:g5:heat-and-insulation:heat}{ऊष्मा} देता है, इसलिए सतह का पानी
सबसे पहले \qty{0}{\celsius} पर पहुँचता है और उसकी बर्फ़ वहीं बनती है
--- और वह \omterm{def:g1:floating-and-sinking:words}{तैरती} है, क्योंकि बराबर जगह के लिए वह पानी से हल्की है,
इसलिए ऊपर ही टिकी रहती है। बर्फ़ (और हिम) का यह ढक्कन एक धीमी लेन है:
वह नीचे के पानी को रोधन देता है, जिसकी \omterm{def:g5:heat-and-insulation:heat}{ऊष्मा} अब बहुत धीरे-धीरे ही ऊपर
रिसती है। गहरा पानी पूरे जाड़े जमने के \omterm{def:g3:temperature-thermometers:temperature}{ताप} से ज़रा ऊपर बना रहता है, और
मछलियाँ उस छत के नीचे \omterm{def:g1:floating-and-sinking:words}{तैरती} हैं जो उनके अनोखे, फैलने वाले पानी ने
उनके लिए बनाई है।
\end{solution}

\begin{solution}{pb:g6:water-states:1}
\textbf{1.} \omterm{def:g6:water-states:changes}{वाष्पन} (गरम सतह से वाष्पीकरण)।
\textbf{2.} तब भी \qty{1}{kg} --- \omterm{def:g3:measuring-mass:mass}{द्रव्यमान} अवस्था के बदलावों को झेल
जाता है।
\textbf{3.} \omterm{def:g6:water-states:changes}{संघनन}।
\textbf{4.} \qty{0}{\celsius} पर।
\textbf{5.} लगभग \qty{1.1}{L}।
\textbf{6.} \omterm{def:g3:measuring-mass:mass}{द्रव्यमान} \qty{1}{kg}; \omterm{def:g6:mass-and-volume:volume}{आयतन} फिर से \qty{1}{L}।
\textbf{7.} $1000 - 250 = \qty{750}{g}$ धारा में आगे चलते हैं।
\textbf{8.} नहीं: वे \qty{250}{g} वाष्प बनकर \omterm{def:g2:air-around-us:air}{हवा} पर सवार हैं,
सही-सलामत; उनका सबसे संभव अगला दरवाज़ा \omterm{def:g6:water-states:changes}{संघनन} है --- पहले बादल, फिर
बारिश।
\textbf{9.} $750 + 250 = \qty{1000}{g}$: पूरा \omterm{def:g3:measuring-mass:units}{किलोग्राम} घर लौट आया।
\textbf{10.} दो बार अनदेखा: समुद्र छोड़ने के ठीक बाद वाष्प के रूप में,
और झरने के कुंड के ऊपर फिर से वाष्प के रूप में --- दोनों बार \omterm{def:g3:solids-liquids-gases:gas}{गैस}
अवस्था में।
\textbf{11.} फूला होना \emph{\omterm{def:g6:mass-and-volume:volume}{आयतन}} की बात है: हिम बर्फ़ के रवों में
\omterm{def:g2:air-around-us:air}{हवा} गुँथी होने से बनता है और ख़ूब जगह घेरता है। हमारे लेबल लगे पानी का
\emph{\omterm{def:g3:measuring-mass:mass}{द्रव्यमान}} ठीक \qty{1}{kg} बना रहा --- जगह के हिसाब से हल्का
होना पदार्थ के हिसाब से हल्का होना नहीं है: सही शब्द-युग्म है
\omterm{def:g3:measuring-mass:mass}{द्रव्यमान} और \omterm{def:g6:mass-and-volume:volume}{आयतन}।
\textbf{12.} \omterm{def:g6:water-states:changes}{वाष्पन}, \omterm{def:g6:water-states:changes}{संघनन}, \omterm{def:g6:water-states:changes}{ठोसीकरण}, \omterm{def:g6:water-states:changes}{गलन} --- और \omterm{def:g2:ice-water-steam:evapcond}{वाष्पित} हुए चौथाई के
लिए: \omterm{def:g6:water-states:changes}{वाष्पन}, \omterm{def:g6:water-states:changes}{संघनन} (बारिश), और फिर \omterm{def:g3:solids-liquids-gases:liquid}{द्रव} के रूप में वापस मिल जाना। चार
दरवाज़े काम आए; \omterm{def:g6:water-states:changes}{ऊर्ध्वपातन} और गैस-से-ठोस वाले दरवाज़े की इस यात्रा में
ज़रूरत नहीं पड़ी।
\end{solution}
